\documentclass[11pt]{article}
\usepackage[top=.76in,bottom=.76in,left=.85in,right=.85in]{geometry}
\usepackage{amsmath,amssymb,amsthm,booktabs,microtype,xcolor}
\usepackage[colorlinks=true,urlcolor=blue!45!black]{hyperref}
\newtheorem{theorem}{Theorem}
\title{Five-step spanning-tree divisibility\\in rank-two graph towers}
\author{}
\date{26 November 2026}
\begin{document}
\maketitle
\vspace{-2em}
\begin{abstract}
The ProveIt upper bound for odd step counts is attained for five steps
at $p=11$, at every layer, by one simple Cayley graph tower.
We give an unconditional proof, an independent determinant check,
and exhaustive maxima at all 22 primes between 11 and 101.
\end{abstract}
\section{Sharpness at eleven}
Let $p>10$ be prime and let $S=(s_1,\ldots,s_5)$ be a list of integer
vectors spanning $\mathbb F_p^2$ after reduction. Repetitions are
allowed. At level $n\ge1$, each vertex $v$ of $(\mathbb Z/p^n\mathbb Z)^2$
and each index $i$ supply an undirected edge from $v$ to $v+s_i$.
Write $\kappa_n(S)$ for the spanning-tree count.
\begin{theorem}
For every $n\ge1$,
\[
\max_S v_{11}(\kappa_n(S))=72(11^n-1)/10-2n.
\]
The list $((1,0),(0,1),(0,2),(1,1),(3,7))$ attains all these maxima.
Its graphs are simple and ten-regular.
\end{theorem}
\begin{proof}
Put $M_j(u)=\sum_i(u\cdot s_i)^{2j}$ over $\mathbb F_p$ and let $r(L)$
be its first nonzero index in projective direction $L$. Newton's
identities give $r(L)\le5$, since otherwise the five squared
projections would all be zero, contrary to spanning. The character
calculation below gives
\[
 A(S)=2\sum_L r(L),\qquad v_p(\kappa_n(S))=A(S)(p^n-1)/(p-1)-2n.
\]
We show $A(S)\le4(p+1)+24$. If $M_1$ is nonzero, its at most two
projective roots give $A(S)\le2(p+1)+16$. Otherwise $M_2$ is nonzero.
Indeed, discard zero steps and group the linear forms by projective
class. Fourth powers of up to five distinct forms are independent:
choose finite slopes and use the first Vandermonde rows and the
invertible degree-four binomial coefficients. If $M_2=0$, every
group would have at least two terms, so there would be exactly two
groups by spanning and the step limit. Independence of the two
squared forms then forces both the first and the second power sums
of squared scales to vanish in each group. By Newton's identities
this is impossible for two nonzero scales, so the groups would
require at least six steps. Thus the quartic $M_2$ is nonzero.
At most four directions have depth above two, proving the bound.

For the stated list the integer coefficients of $M_1$ are $11,44,55$.
In directions $(1,t)$ for $0\le t\le10$, then $(0,1)$, the $M_2$
residues are
\[
 6,\ 2,\ 0,\ 0,\ 0,\ 7,\ 4,\ 5,\ 0,\ 10,\ 10,\ 10.
\]
At each of its four zeros, $(M_3,M_4,M_5)=(0,0,5)$ modulo eleven.
Thus $A=2(8\cdot2+4\cdot5)=72$. The steps are nonzero and distinct up
to sign modulo eleven, so no loops or parallel edges occur.
\end{proof}
\section{The character calculation}
At conductor $p^h$ let $u$ be a primitive covector. The Laplacian
eigenvalue is $\sum_i(2-Z^{a_i}-Z^{-a_i})$, where $a_i=u\cdot s_i$
and $Z=\zeta_{p^h}$. Write $T=Z-1$ and $e=(p-1)p^{h-1}$. The
cyclotomic valuation gives $v_p(T)=1/e$. Modulo $p$, the invertible
substitution $1+T=\exp(z)$ through degree ten yields coefficients
$-2M_j(\bar u)/(2j)!$ at $z^{2j}$ and zero at odd degrees.
Thus the first unit coefficient of the integral $T$-expansion
is at degree $2r(L)\le10$. All earlier nonconstant terms have
valuation at least $1+1/e$, and the constant term is exactly zero.
Later terms have valuation strictly larger than $2r(L)/e$.
This proves the eigenvalue valuation even when $2r(L)=e=10$.
Each residue direction has $p^{h-1}$ projective lifts, each giving
$e$ primitive characters. Summation over $h$ and division by
$p^{2n}$ in the matrix-tree theorem give the formula above.

Independently, fraction-free elimination on the $120\times120$
reduced integer Laplacian for the witness at $n=1$ gives valuation
70 and quotient by $11^{70}$ congruent to 1 modulo eleven.
\section{Other primes}
\paragraph{Congruence obstruction.}
For $m$ steps and $p>2m$, a direction of depth $m$ has squared
projections with monic polynomial $z^m-c$, $c\ne0$. Its $m$
distinct roots are nonzero squares. Their ratios to any one root
form the group of $m$th roots of unity, so $2m\mid p-1$.
If this fails, all depths are at most $m-1$. For five steps the
upper bound drops to $4(p+1)+16$. In general, by the
\href{https://github.com/VladimirReshetnikov/ProveIt/blob/21c0252ffe8611baa0fbf528fbb003d6c42ca9af/Analysis/GeneralizedRiemannHypothesis/Research/grh_transfer_extensions_2026-11-26/sections/graph_extremum.tex}{ProveIt first-moment bound},
the first nonzero moment polynomial has index $a\le k=\lfloor m/2\rfloor$.
It has at most $2a$ exceptional directions, so the improved bound
is $2a(p+1+2m-2-2a)$, increasing in $a$ up to $k$.

\paragraph{Infinite family.}
For $p\equiv1\pmod{10}$, use $(\xi^i,\xi^{2i})$ for $0\le i<5$,
where $\xi$ has order five. The geometric sum in the monomial
of degree $2j$ and $Y$-degree $b$ is supported on $5\mid 2j+b$.
Hence $M_1=0$, $M_2=20X^3Y$, and both axis directions have
depth five. The other directions have depth two, proving the
lower bound $4(p+1)+12$ on that progression.

\paragraph{Exhaustive search.}
For $M_1=0$, a change of basis sends two independent steps to
$(1,0),(0,1)$, without changing $A$. Use sign representatives in
the remaining triple, optionally including zero. The key
$K(x,y)=(x^2,xy,y^2)$ is injective on these representatives.
For each sorted pair $u\le v$ in the triple the remaining $w\ge v$
must have key $(-1,0,-1)-K(u)-K(v)$. This exhausts all normalized
lists. In every tested case the best score exceeds the
$M_1\ne0$ ceiling $2(p+1)+16$.
The table gives $A_{\max}$, and the CSV also records the counts, with and without
zero residues. Allowing zero changes eight of the 22 maxima.
An integer step divisible by $p$ in both coordinates counts as
a zero residue even if it is nonzero at higher layers.

\begin{center}
\small
\begin{tabular}{r rr @{\qquad} r rr @{\qquad} r rr}
\toprule
$p$ & Nonzero & Zero OK & $p$ & Nonzero & Zero OK & $p$ & Nonzero & Zero OK\\
\midrule
11 & 72 & 72 & 41 & 180 & 184 & 73 & 306 & 312\\
13 & 58 & 64 & 43 & 180 & 184 & 79 & 328 & 328\\
17 & 80 & 88 & 47 & 200 & 200 & 83 & 344 & 344\\
19 & 84 & 88 & 53 & 224 & 224 & 89 & 368 & 376\\
23 & 96 & 96 & 59 & 248 & 248 & 97 & 402 & 408\\
29 & 128 & 128 & 61 & 260 & 260 & 101 & 424 & 424\\
31 & 148 & 148 & 67 & 280 & 280 &  &  & \\
37 & 160 & 160 & 71 & 300 & 300 &  &  & \\
\bottomrule
\end{tabular}
\end{center}
\paragraph{Scope.}
This resolves one sharpness case of the odd-step question in the
\href{https://github.com/VladimirReshetnikov/ProveIt/blob/21c0252ffe8611baa0fbf528fbb003d6c42ca9af/Analysis/GeneralizedRiemannHypothesis/Research/grh_transfers_lubin_tate_and_graph_towers/graphs.tex}{ProveIt graph-tower report}.
The transfer is of cyclotomic specialization, not of an RH or GRH
theorem. For general graph-tower growth see
\href{https://doi.org/10.5802/alco.304}{On $\mathbb Z_\ell^d$-towers of graphs}.
The exact maximum at arbitrary odd step counts and primes remains
open here; no claim of historical priority is made.
\end{document}
